\begin{PersianAbstract}
	
	بیماری دیستروفی عضلانی دوشن 
	\footnote{
		\lr{Duchenne Muscular Dystrophy (DMD)}
	}
	یکی از شایع‌ترین و شدیدترین اختلالات ژنتیکی عضلانی در کودکان است که با ضعف پیشرونده عضلات اسکلتی همراه بوده و به‌مرور زمان منجر به کاهش توانایی‌های حرکتی می‌شود. پایش دقیق و کمی عملکرد حرکتی بیماران مبتلا به این بیماری، نقش مهمی در ارزیابی روند پیشرفت بیماری، بررسی اثربخشی درمان‌ها و طراحی برنامه‌های توان‌بخشی ایفا می‌کند. در حال حاضر، یکی از ابزارهای رایج برای ارزیابی عملکرد حرکتی این بیماران، مقیاس
	 \lr{NSAA}
	\footnote{ \lr{North Star Ambulatory Assessment} }
	است که بر اساس مشاهده بالینی و قضاوت درمانگر، نمره‌ای گسسته بین صفر تا دو به هر حرکت اختصاص می‌دهد. با وجود کاربرد گسترده این مقیاس، ماهیت کیفی و وابستگی آن به ارزیابی انسانی می‌تواند منجر به بروز خطاهای بین ارزیابان و کاهش دقت تکرارپذیری نتایج شود.
	
	هدف این پژوهش، طراحی و توسعه یک سامانه مبتنی بر حسگرهای اینرسیایی به‌منظور ارزیابی کمی شاخص‌های عملکرد موتوری در بیماران مبتلا به دوشن و نزدیک‌سازی فرآیند نمره‌دهی به یک چارچوب داده‌محور است. در این راستا، مجموعه‌ای از حسگرهای اینرسیایی شامل شتاب‌سنج و ژیروسکوپ بر روی بخش‌های کلیدی بدن نصب شده و داده‌های شتاب خطی و سرعت زاویه‌ای حین اجرای حرکات منتخب 
	\lr{NSAA} 
	ثبت گردید. داده‌های خام جمع‌آوری‌شده پس از انتقال به محیط متلب
	 \lr{MATLAB}،
	 تحت فرآیند پیش‌پردازش شامل حذف نویز، فیلترگذاری دیجیتال و در صورت نیازی نرمال‌سازی زمانی قرار گرفتند. در ادامه، با استخراج ویژگی‌های مناسب از حوزه زمان، مجموعه داده‌ای ساختاریافته برای تحلیل‌های بعدی آماده شد.
	
	به‌منظور پیش‌بینی نمره هر حرکت، از الگوریتم‌های یادگیری ماشین استفاده گردید و نتایج حاصل با ارزیابی بالینی مقایسه شد. در فاز مقدماتی این پژوهش، بازدید میدانی و مشاهده وضعیت حرکتی حدود ۱۲ بیمار مبتلا به دیستروفی عضلانی دوشن با همکاری تیم بالینی انجام شد که درک عمیق‌تری از الگوهای حرکتی و محدودیت‌های عملی سیستم فراهم ساخت. نتایج نشان داد که رویکرد مبتنی بر حسگرهای اینرسیایی و تحلیل داده‌محور می‌تواند به‌عنوان ابزاری مکمل در کنار ارزیابی بالینی، به افزایش دقت، عینیت و قابلیت تکرارپذیری فرآیند نمره‌دهی کمک نماید.
	
	این سامانه می‌تواند در آینده به‌عنوان زیرساختی برای پایش طولی بیماران، توسعه سامانه‌های ارزیابی از راه دور و بهبود تصمیم‌گیری‌های درمانی مورد استفاده قرار گیرد.
	
 \keywords{
 	دیستروفی عضلانی دوشن، حسگر اینرسیایی، تحلیل حرکت، یادگیری ماشین، مقیاس 
 	\lr{NSAA}
 	، ارزیابی عملکرد حرکتی
 	}
	
\end{PersianAbstract}
